\section*{Capítulo \ref{ch:b1:digestion-absorption} --- Digestão e absorção}

\begin{solution}{exo:b1:digestion-absorption:1}
Boca: mastigação, amilase salivar. Esôfago: transporte. Estômago:
ácido, pepsina, armazenamento e mistura. Intestino delgado: a bile e
as \omterm{def:b1:enzymes:enzyme}{enzimas} pancreáticas completam a \omterm{def:b1:digestion-absorption:nutrition}{digestão}; absorção. Intestino
grosso: água e sal recuperados, \omterm{def:b1:respiration-fermentation:fermentation}{fermentação} pelos micróbios. Reto:
armazenamento e egestão.
\end{solution}

\begin{solution}{exo:b1:digestion-absorption:2}
Pepsina no estômago a pH 2; tripsina, quimotripsina e
carboxipeptidases do pâncreas no intestino delgado a pH 8; peptidases
da borda em escova a pH 7. Como zimogênios, para que não digiram as
\omterm{def:b1:cell-unit-of-life:cell}{células} que as fabricam; só são ativadas no lúmen.
\end{solution}

\begin{solution}{exo:b1:digestion-absorption:3}
Glicose: simporte com sódio para dentro, carreador para fora, rumo à
veia porta e ao fígado. \omterm{def:b1:proteins:aminoacid}{Aminoácido}: simporte com sódio para dentro,
carreador para fora, veia porta. \omterm{def:b1:lipids:fattyacid}{Ácido graxo}: difusão para dentro a
partir de uma micela, reesterificado, empacotado num quilomícron,
exocitado no quilífero e na \omterm{def:b1:mammal-organization:compartments}{linfa}.
\end{solution}

\begin{solution}{exo:b1:digestion-absorption:4}
Os sais biliares emulsionam a \omterm{def:b1:lipids:triglyceride}{gordura} em gotículas de um micrômetro e
levam os produtos da lipase em micelas até os \omterm{def:b1:digestion-absorption:surface}{enterócitos}. Sem eles a
\omterm{def:b1:lipids:triglyceride}{gordura} permanece em gotas grandes com quase nenhuma superfície para a
lipase, e os \omterm{def:b1:lipids:fattyacid}{ácidos graxos} liberados não podem ser entregues à
membrana: a \omterm{def:b1:lipids:triglyceride}{gordura} atravessa sem ser digerida.
\end{solution}

\begin{solution}{exo:b1:digestion-absorption:5}
$100/162 = \qty{0.62}{mol}$ de unidades de glicose, portanto cerca de
\qty{0.62}{mol} de ligações e \qty{0.62}{mol} de água (\qty{11}{g});
$620/180 = \qty{3.4}{mmol/min}$ de glicose absorvida.
\end{solution}

\begin{solution}{exo:b1:digestion-absorption:6}
O ácido libera a secretina, que faz o pâncreas secretar bicarbonato,
que neutraliza o quimo; a \omterm{def:b1:lipids:triglyceride}{gordura} e os peptídeos liberam a
colecistocinina, que contrai a vesícula biliar (a bile entra),
estimula as \omterm{def:b1:enzymes:enzyme}{enzimas} pancreáticas e inibe o esvaziamento gástrico, de
modo que o duodeno recebe quimo apenas na velocidade com que consegue
neutralizá-lo e digeri-lo.
\end{solution}

\begin{solution}{exo:b1:digestion-absorption:7}
A glicose foi concentrada oito vezes contra o seu gradiente:
\omterm{def:b1:membranes-transport:transporters}{transporte ativo}. Sem sódio o transporte para no equilíbrio: a
captação é acoplada ao gradiente de sódio (o \omterm{prop:b1:membranes-transport:secondary}{simportador}
sódio--glicose), e não diretamente ao \omterm{def:b1:water-small-molecules:atp}{ATP}.
\end{solution}

\begin{solution}{exo:b1:digestion-absorption:8}
Os quilomícrons, com um micrômetro de diâmetro, não conseguem
atravessar a parede fechada do capilar, mas podem entrar nos
quilíferos, abertos na ponta; a \omterm{def:b1:mammal-organization:compartments}{linfa} junta-se ao sangue no pescoço, e
por isso a \omterm{def:b1:lipids:triglyceride}{gordura} alcança o corpo inteiro antes do fígado, que de
outro modo a captaria toda primeiro. Entrando diretamente no sangue
portal, os quilomícrons entupiriam os sinusoides do fígado e
saturariam sua capacidade de processar \omterm{def:b1:lipids:triglyceride}{gordura} depois de cada refeição.
\end{solution}

\begin{solution}{exo:b1:digestion-absorption:9}
Ela umedece e faz flutuar a fibra para a ruminação; seu bicarbonato e
seu fosfato tamponam os ácidos que os micróbios produzem, mantendo o
rúmen perto de pH 6,5; e ela transporta ureia vinda do sangue, que os
micróbios convertem em amônia e depois na sua própria \omterm{def:b1:proteins:peptide}{proteína},
reciclando nitrogênio que de outro modo se perderia na urina.
\end{solution}

\begin{solution}{exo:b1:digestion-absorption:10}
Sem a amilase, as proteases, a lipase e o bicarbonato pancreáticos: o
\omterm{def:b1:carbohydrates:polysaccharide}{amido} é digerido em parte apenas pela saliva e pela borda em escova;
as \omterm{def:b1:proteins:peptide}{proteínas} são cortadas pela pepsina mas não até um tamanho
absorvível; a \omterm{def:b1:lipids:triglyceride}{gordura} não é digerida de modo nenhum e aparece em fezes
volumosas, pálidas e gordurosas. Açúcares e \omterm{def:b1:carbohydrates:glycosidic}{dissacarídeos} (as \omterm{def:b1:enzymes:enzyme}{enzimas}
da borda em escova permanecem), parte do \omterm{def:b1:carbohydrates:polysaccharide}{amido} e pequenos peptídeos
ainda podem ser aproveitados; a \omterm{def:b1:lipids:triglyceride}{gordura} e a maior parte da \omterm{def:b1:proteins:peptide}{proteína}
não.
\end{solution}

\begin{solution}{exo:b1:digestion-absorption:11}
As fezes cecais carregam a \omterm{def:b1:proteins:peptide}{proteína} microbiana e as vitaminas do
complexo B produzidas no ceco, que fica além do intestino delgado e
por isso não podem ser digeridas na primeira passagem; comidas, elas
atravessam o estômago e o intestino delgado e são absorvidas. A vaca
fermenta antes do estômago, de modo que seus micróbios são digeridos
adiante, sem qualquer segunda passagem. O ceco do cavalo também fica
adiante, e o cavalo nem come suas fezes nem dispõe de outra via: sua
\omterm{def:b1:proteins:peptide}{proteína} microbiana se perde.
\end{solution}

\begin{solution}{exo:b1:digestion-absorption:12}
Os micróbios dão à vaca a energia da \omterm{def:b1:carbohydrates:polysaccharide}{celulose} sob a forma de \omterm{def:b1:lipids:fattyacid}{ácidos
graxos}, \omterm{def:b1:proteins:peptide}{proteína} feita de capim pobre e até de ureia, e todas as
vitaminas do complexo B; a vaca lhes dá um tonel quente, tamponado e
anaeróbio, mantido cheio e mexido, e cem litros de saliva por dia. Os
custos do projeto: um décimo da energia sai como metano, nenhuma
glicose é absorvida e por isso o fígado tem de fabricá-la toda a
partir do propionato, a \omterm{def:b1:digestion-absorption:nutrition}{digestão} leva três dias, e o tonel e seu
conteúdo pesam um quinto do animal --- uma vaca não consegue disparar,
e um leão consegue.
\end{solution}

\begin{solution}{pb:b1:digestion-absorption:1}
\textbf{1.} Fibra \qty{5}{kg}, \omterm{def:b1:carbohydrates:polysaccharide}{amido} e açúcares \qty{2}{kg}, \omterm{def:b1:proteins:peptide}{proteína}
\qty{1.5}{kg}.
\textbf{2.} $0.7\times 5 + 2 = \qty{5.5}{kg}$.
\textbf{3.} $5500\times 17.5 = \qty{96}{MJ}$; AGV $0.75\times 96 =
\qty{72}{MJ}$.
\textbf{4.} $0.08\times 96 = \qty{7.7}{MJ}$; $7700/55 = \qty{140}{g}$
de metano, cerca de \qty{200}{L}.
\textbf{5.} $0.15\times 5500 = \qty{825}{g}$.
\textbf{6.} A maior parte da \omterm{def:b1:proteins:peptide}{proteína} alimentar é degradada pelos
micróbios no rúmen e reconstruída como \omterm{def:b1:proteins:peptide}{proteína} microbiana; os
micróbios seguem para o abomaso e o intestino delgado, onde são
digeridos como qualquer carne.
\textbf{7.} \qty{30}{\%} da fibra, \qty{1.5}{kg}: fermentada em parte
no cólon, na maior parte excretada.
\textbf{8.} \omterm{def:b1:proteins:peptide}{Proteína} alimentar que escapa $0.4\times 1500 = \qty{600}{g}$,
digerida \qty{360}{g}; microbiana \qty{825}{g}; total \qty{1185}{g}
$\times 17.5 = \qty{20.7}{MJ}$.
\textbf{9.} $72 + 20.7 = \qty{93}{MJ}$.
\textbf{10.} $72/93 = \qty{77}{\%}$ (cerca de \qty{60}{\%} quando o
calor da \omterm{def:b1:respiration-fermentation:fermentation}{fermentação} e a energia da contribuição do cólon entram em
balanços mais completos; nesta conta, três quartos).
\textbf{11.} \qty{93}{MJ} fica aquém de \qty{110}{}: a vaca tem de
comer mais (\qty{12}{kg}), ou capim melhor, ou recorrer à sua \omterm{def:b1:lipids:triglyceride}{gordura}.
\textbf{12.} Propionato; \omterm{def:b1:biosyntheses-integration:gluconeogenesis}{gliconeogênese} no fígado.
\textbf{13.} $0.9\times 2000\times 17.5 = \qty{31.5}{MJ}$.
\textbf{14.} $0.7\times 1500\times 17.5 = \qty{18.4}{MJ}$.
\textbf{15.} $0.45\times 5 = \qty{2.25}{kg}$; AGV $0.75\times 2250\times
17.5 = \qty{29.5}{MJ}$.
\textbf{16.} $0.15\times 2250 = \qty{340}{g}$, perdidos nas fezes (o
intestino posterior fica além do intestino delgado).
\textbf{17.} $31.5 + 18.4 + 29.5 = \qty{79}{MJ}$; parcela dos AGV $29.5/79 =
\qty{37}{\%}$.
\textbf{18.} Não: precisa de $110/79\times 10 = \qty{14}{kg}$ de capim,
e consegue comendo durante dezesseis horas por dia e fazendo o
alimento passar duas vezes mais depressa.
\textbf{19.} Vaca $93/10 = \qty{9.3}{MJ/kg}$; cavalo $79/10 =
\qty{7.9}{MJ/kg}$.
\textbf{20.} Ingestão bruta $10\,000\times 0.85\times 17.5 =
\qty{149}{MJ}$ (excluídas a lignina e as cinzas); o metano,
\qty{7.7}{MJ}, é \qty{5}{\%} disso; cem vacas: $100\times 140\times 365 =
\qty{5.1}{t}$ de metano por ano.
\textbf{21.} Vaca $10\times 70/24 = \qty{29}{kg}$ de matéria seca no
intestino (mais dez vezes isso de água); cavalo $10\times 36/24 =
\qty{15}{kg}$.
\textbf{22.} Em princípio os \qty{340}{g} inteiros, que valem
\qty{6}{MJ} e seus \omterm{def:b1:proteins:aminoacid}{aminoácidos} essenciais; o conteúdo cecal de um
cavalo não é separado numa fração mole e rica em nutrientes como o de
um coelho, e um animal de meia tonelada não consegue recolher do chão
as próprias fezes em estado aproveitável.
\textbf{23.} Os micróbios do rúmen fabricam \omterm{def:b1:proteins:peptide}{proteína} a partir do
nitrogênio da ureia e do carbono da palha, e a vaca digere os
micróbios: ela quase não precisa de \omterm{def:b1:proteins:peptide}{proteína} alimentar. Os micróbios
do cavalo ficam além do seu intestino delgado, de modo que a \omterm{def:b1:proteins:peptide}{proteína}
deles se perde, e o cavalo tem de encontrar seus \omterm{def:b1:proteins:aminoacid}{aminoácidos} no
próprio alimento.
\textbf{24.} O cavalo absorve glicose e armazena \omterm{def:b1:carbohydrates:polysaccharide}{glicogênio} nos
músculos para disparar em anaerobiose; a vaca não absorve nenhuma,
fabrica glicose lentamente a partir do propionato e carrega um tonel
de cem litros --- feita para pastar, não para correr.
\textbf{25.} Cerca de três quartos da energia metabolizável da vaca
neste balanço (\qtyrange{60}{70}{\%} numa contabilidade completa) como
\omterm{def:b1:lipids:fattyacid}{ácidos graxos} voláteis, contra um terço no cavalo; \qty{9.3}{MJ}
contra \qty{7.9}{MJ} por quilograma de capim.
\end{solution}
